\documentclass{situationist-preprint}

\title{Russellian Monism Completed}
\msid{MND-01-c}
\msfamily{MND-01}
\msdiscipline{Philosophy of Mind}
\mstype{Alternative proof}
\mskeywords{Russellian monism; protophenomenal properties; structuralism; intrinsic natures; physicalism; grounding; conceivability}

\begin{document}
\maketitle

\begin{abstract}
MND-01-a reaches Russellian monism through the two-dimensional argument: given materialism, the premises that link the conceivability of zombies to their primary possibility, and primary possibility to secondary possibility or Russellian monism, leave Russellian monism as the only option. Those premises, and above all the step from ideal conceivability to primary possibility, are the most disputed in the literature. We reach the conclusion of MND-01 without them. From structuralism about physics, the principle that relations require relata with intrinsic natures, the non-structural character of phenomenal properties, the physicalism of MND-02 and a closure principle for structural grounding, we derive that every phenomenal property is grounded in non-structural, non-mental fundamental natures together with structure: protophenomenal Russellian monism. The explanatory gap then follows as a gap between modes of presentation. The route completes the structural picture of physics rather than arguing from what can be conceived.
\end{abstract}

\section{The premise avoided}\label{s:avoid}

The principal manuscript of MND-01 argues that the hard problem of consciousness can be solved: consciousness is the intrinsic nature of physical structure, and the explanatory gap arises from two modes of presentation of the same facts. Its main argument is the two-dimensional argument, stated as Assumption~4.1 of MND-01-a. With $P$ the conjunction of microphysical truths and $Q$ a phenomenal truth, the assumption has four clauses: (a) $P \wedge \neg Q$ is ideally conceivable; (b) if it is ideally conceivable, it is primarily possible; (c) if it is primarily possible, then it is secondarily possible or Russellian monism is true; (d) if it is secondarily possible, materialism is false. From these, Theorem~4.2 of MND-01-a derives that materialism is false or Russellian monism is true, and Corollary~4.3 that, given materialism, Russellian monism is true.

Clause (b) is the most disputed premise in this argument. Type-B physicalists reject it, holding that the zombie hypothesis is ideally conceivable without being primarily possible \citep{Loar1990,Papineau2002}; others reject the two-dimensional framework that gives it content \citep{BlockStalnaker1999}. Clause (a) is rejected by illusionists. Clause (c) is plausible only given a specific theory of the semantics of physical terms. A Russellian monism that rests on all four is exposed to every objection against any of them.

The present manuscript reaches the conclusion of MND-01 without Assumption~4.1. We do not use conceivability at all, ideal or otherwise, and we make no assumption linking any epistemic modality to any metaphysical one. What we keep from MND-01-a is its structuralism about physics (MND-01-a, P2) and its materialism, which is the physicalism of MND-02. The work Assumption~4.1 does in the principal argument is done here by three metaphysical premises: that relations require relata with natures that are not themselves relational, that phenomenal properties are not structural, and that what is wholly grounded in structure alone is structural. From these and physicalism we derive not only that Russellian monism is true but which version of it: the protophenomenal version, which MND-01-a reaches only with the further help of MND-18.

The title reflects the form of the argument. Physics, by MND-01-a's P2, gives an incomplete description of the physical world: it says what plays which role and leaves open what the role-players are in themselves. The two-dimensional argument shows that, given materialism, the gap must be filled in a Russellian way, by showing that every other way of filling it, or of leaving it empty, leads to the falsity of materialism. The argument given here fills the gap directly. It asks what the physical world must contain if structuralism is true, if there are phenomenal properties, and if everything is physical, and shows that the answer is a completion of the structural picture by intrinsic natures that ground experience. Russellian monism is that completion.

Section~\ref{s:history} reviews the development of Russellian monism and the objections it has faced. Section~\ref{s:frame} states the framework and Section~\ref{s:arg} the argument. Section~\ref{s:causal} shows that the completed picture gives experience a causal role. Section~\ref{s:compare} compares the two routes. Section~\ref{s:obj} considers objections, and Section~\ref{s:not} says what the argument does not establish.

\section{Russellian monism: a short history}\label{s:history}

\subsection{Russell and Eddington}

In \emph{The Analysis of Matter}, \citet{Russell1927} argued that physics gives us knowledge of the abstract structure of the physical world and no knowledge of its intrinsic character. The theoretical entities of physics are known only through their relations to one another and to our percepts, and a description of relations is a description of structure. Russell drew from this a conclusion about the mind: since the one part of the world whose intrinsic character we know is our own experience, and since our brains are part of the physical world, the intrinsic character of at least some physical events is known to us, and it is experiential. \citet{Eddington1928} put the same thought vividly, describing physics as a scheme of symbols and pointer readings that tells us nothing of the background those symbols stand for, and suggesting that in our own minds we encounter that background directly.

\citet{Newman1928} raised an objection to Russell's structuralism that bears on our argument. If all we know of the physical world is its structure, then what we know is nearly trivial: any collection of the right cardinality can be organized to have any given structure. Russell accepted the point. Its lesson, which \citet{DemopoulosFriedman1985} drew out, is that a purely structural description underdetermines its subject matter except as to cardinality, and that whatever more we know about the physical world must come from somewhere other than the structure itself. MND-10 applies the same point to quality spaces: a relational description of a quality space determines which quality occupies which place only up to automorphism and never determines the intrinsic character of an occupant (MND-10-a, Lemma~5.1 and Proposition~6.1).

\subsection{The identity theorists}

The view was revived by philosophers who saw in it a way to make the mind--brain identity theory intelligible. \citet{Feigl1958} argued that the referents of physical terms are known only by description, as whatever plays certain roles, and that the identity of mental with physical states is the discovery that what plays the neural roles is what we are acquainted with in experience. \citet{Maxwell1978} defended a similar position against Kripke's modal argument, arguing that physical terms are not rigid designators of the properties that occupy their roles and that the apparent contingency of psychophysical identities is therefore explicable. \citet{Lockwood1989} developed the view in connection with the interpretation of quantum mechanics. In each case the motivation was the same: if physics leaves the intrinsic natures of its entities open, the mind can be the intrinsic nature of the brain without being anything over and above the physical.

\subsection{The contemporary revival}

\citet{Chalmers1996} gave the view its current name and its place in the dialectic: Russellian monism is the one position that respects both the anti-materialist arguments and the causal closure of the physical, since the phenomenal or protophenomenal properties it posits are the categorical bases of the dispositions physics describes and so are causally efficacious without overdetermination. \citet{Stoljar2001} drew the distinction between theory-based and object-based conceptions of the physical, on which Russellian monism is physicalist in the second sense and not in the first. \citet{Strawson2006} argued that any realistic physicalism must be panpsychist, since experience cannot emerge from the wholly non-experiential. \citet{Pereboom2011} argued that Russellian monism is a live option and that the proper attitude to it is one of open-mindedness about the intrinsic natures of matter. \citet{AlterNagasawa2012} gave a precise formulation of the view, distinguishing panpsychist, panprotopsychist and physicalist versions, and the collection edited by \citet{AlterNagasawa2015} surveys the state of the debate. \citet{Goff2017} gave an extended defence of a panpsychist version, arguing for it from the conjunction of anti-physicalist arguments with a demand for a parsimonious and unified theory.

\subsection{The objections}

Three objections have shaped the debate. The first is the \emph{combination problem}: how the phenomenal or protophenomenal natures of microphysical entities combine to ground the unified experience of a subject \citep{Coleman2014,Chalmers2017}. MND-18 argues that the problem cannot be solved for panpsychism, since the experiences of micro-subjects cannot ground the experience of a further subject (MND-18-a, Theorem~3.2), and that the Russellian monism of MND-01 must therefore be protophenomenal; MND-01-d argues that the residual boundary and palette problems are problems of articulating a grounding relation, not of showing that it can exist (MND-01-d, Proposition~2.1).

The second is the \emph{objection from pure powers}. If fundamental properties are pure powers, wholly constituted by their dispositional roles, there are no categorical natures for experience to be \citep{Bird2007}. The third is the \emph{objection from humility}. If physics tells us only about roles, we are ignorant of the intrinsic natures of physical things, and Russellian monism claims to know too much about them \citep{Langton1998,Lewis2009}. We take up the second and third in Section~\ref{s:obj}. The first concerns the articulation of the view rather than the argument for it, and is outside our scope.

\subsection{Arguments for the view}

The arguments for Russellian monism in the literature are of three kinds. The first is the two-dimensional argument, which MND-01-a adopts: given the conceivability of zombies and the links between conceivability and possibility, Russellian monism is the only form of materialism left standing \citep{Chalmers2009}. The second is the argument from mental causation: Russellian monism is the only view on which experience is both non-structural and causally efficacious in a causally closed physical world \citep{Chalmers1996}. The third is the argument from the need for intrinsic natures: physics describes relations, relations need relata, and the relata must have intrinsic natures, of which experience is the only candidate we know \citep{Russell1927,Strawson2006}.

The present argument is of the third kind, but in a form that avoids the step usually made at its end: the claim that experience is the \emph{only} candidate for the intrinsic natures, which leads to panpsychism. Our argument does not claim that the intrinsic natures of matter are experiences, or even that experience is the only kind of intrinsic nature we can conceive. It claims that the grounds of experience must include intrinsic natures that are physical and not mental, and that this is what protophenomenal Russellian monism says.

\section{Framework}\label{s:frame}

\begin{definition}[Structural property]\label{d:struct}
A property is \emph{structural} if its instantiation consists in its bearer's occupying a place in a pattern of nomic, causal, spatiotemporal or mathematical relations: what it is for a thing to have the property is for it to stand in such relations to other things, or to be disposed to.
\end{definition}

Mass, charge and spin, as physics characterizes them, are structural in this sense: to have unit charge is to be disposed to produce and respond to electromagnetic fields in a lawfully specified way. So are functional and dispositional properties at higher levels. Definition~\ref{d:struct} concerns what instantiation of the property \emph{consists in}, not merely how the property is described: a property that is described relationally but whose instantiation consists in something else, as the property of being the cause of these symptoms may be instantiated by having a certain molecular constitution, is not structural for that reason alone.

\begin{definition}[Intrinsic nature]\label{d:nature}
An \emph{intrinsic nature} of a thing is a property it has that is not structural and that it has independently of its relations to other things.
\end{definition}

\begin{definition}[Protophenomenal Russellian monism]\label{d:prm}
\emph{Protophenomenal Russellian monism} is the conjunction of three theses: (i) the fundamental properties that physics characterizes are structural, and the fundamental entities that have them have intrinsic natures; (ii) these intrinsic natures are protophenomenal: they are not experiences and do not by themselves make their bearers subjects, but can in suitable combination ground experiences; and (iii) the experiences of subjects are wholly grounded in these natures together with the structures in which their bearers stand.
\end{definition}

Definition~\ref{d:prm} is the definition of MND-18-b (Definition~1.1), with the notion of a protophenomenal property taken from MND-18-a (Definition~2.3). It is also the view that Volume~III of MND-01 develops, where the intrinsic natures of fundamental entities are defined as protophenomenal (Definition~21.2) and shown to be physical in the sense of MND-02 (Theorem~22.5).

We now state the premises. The first is retained from MND-01-a.

\begin{principle}[Structuralism]\label{p:struct}
The fundamental properties that physics characterizes are structural: physics specifies the structural and dynamical properties of the world and does not specify the intrinsic natures, if any, of the things that have them.
\end{principle}

\ref{p:struct} is P2 of MND-01-a, stated in the terminology of Definition~\ref{d:struct}. It is a claim about the content of physical theory, supported by the way physics introduces its fundamental magnitudes, through their roles in laws.

\begin{principle}[Relata]\label{p:relata}
Whatever stands in relations has an intrinsic nature. In particular, the entities that occupy the roles described by physics have intrinsic natures.
\end{principle}

\ref{p:relata} is the premise of the third kind of argument above. Its support is a regress argument familiar from \citet{Russell1927} and \citet{Blackburn1990}. If an entity's every property consisted in its standing in relations to other entities, and theirs in their standing in relations to yet others, then nothing would be specified anywhere except a pattern of relations among items that are themselves nothing but places in the pattern. A world of that kind would differ from a mathematical structure only in being called concrete, and Newman's point shows that the structure alone would determine nothing but cardinality. That the physical world is more than an abstract structure of a given cardinality is the content of its being concrete, and \ref{p:relata} records it. Note that \ref{p:relata} does not say that structural properties reduce to intrinsic natures, or that dispositions are grounded in categorical bases; it says only that the bearers of structural properties have, in addition, properties that are not structural. We return to the point in Section~\ref{s:obj}.

\begin{principle}[Non-structurality]\label{p:nonstruct}
Phenomenal properties are not structural, and they are instantiated.
\end{principle}

\ref{p:nonstruct} is supported by MND-10. Qualia exist (MND-10-a, Proposition~4.3), and they are intrinsic to the experiences that have them (MND-10-a, Theorem~5.4). The argument for intrinsicness is an application of Newman's point: a relational description of a quality space is satisfied equally by every structure-preserving permutation of its occupants (MND-10-a, Lemma~5.1), so qualia are not constituted by their places in such a space (MND-10-a, Lemma~5.2). The same argument applies to the wider causal and functional roles of experiences, since a description of those roles is also relational and is also satisfied under permutations of their occupants that preserve the roles. That functionally alike subjects could differ in their qualia, as MND-22 holds for inverted spectra, is a consequence.

\begin{assumption}[Physicalism]\label{a:phys}
Every instantiated property is physical in the sense of MND-02: it is a fundamental property of a foreseeable extension of current physics, or is wholly grounded in such properties. The fundamental properties of a foreseeable extension either play roles in its dynamical laws or are the categorical realizers of properties that do, and none of them is mental.
\end{assumption}

Assumption~\ref{a:phys} is the conjunction of Definitions~2.2, 2.3 and~2.4 of MND-02-a with its Theorem~3.6, which establishes physicalism so defined from four premises about the actual fundamental theory, the most substantive of which is that no actually instantiated fundamental property is mental (MND-02-a, Assumption~3.3). This is the materialism that MND-01-a takes as given in its Corollary~4.3, and it is the one premise that the present argument shares with the principal one apart from \ref{p:struct}. Its support is empirical, from the absence of mental forces in physiology \citep{Papineau2001}, and does not depend on any claim about conceivability.

\begin{principle}[Structural closure]\label{p:closure}
A property wholly grounded in structural properties alone is structural.
\end{principle}

\ref{p:closure} says that structure grounds only structure. Its support is that grounding preserves the kind of fact grounded in the following minimal sense: a property wholly grounded in the instantiation of relations by things, and in nothing else about those things, has instantiation conditions that consist in the instantiation of those relations, and so is structural by Definition~\ref{d:struct}. A molecule's shape, a crystal's lattice symmetry and a computer's functional organization are wholly grounded in the structural properties of their constituents and are themselves structural. \ref{p:closure} is the metaphysical counterpart of the epistemic principle of structural closure stated in MND-01-b (P2 there), and it is weaker than the principle, rejected by \citet{Strawson2006}, that nothing non-experiential can ground experience: it says only that the purely structural cannot ground the non-structural.

\begin{assumption}[Uniformity]\label{a:unif}
Fundamental entities of the same physical kind have intrinsic natures of the same kind, and the entities that compose the brains of conscious subjects are of the same physical kinds as those that compose other matter.
\end{assumption}

Assumption~\ref{a:unif} is a principle of the uniformity of nature applied to intrinsic natures. Electrons in a brain and electrons in a star play the same role; nothing in physics distinguishes them, and a view on which their natures differed in kind would posit a difference that makes no difference to anything physics describes and that would be arbitrary in its distribution. The second clause is a commonplace of chemistry and biology.

\section{The argument}\label{s:arg}

\begin{lemma}\label{l:natures}
The fundamental entities that occupy the roles described by physics have intrinsic natures, and these natures are physical: they are fundamental properties of the actual fundamental theory, as categorical realizers of properties that play roles in its dynamical laws.
\end{lemma}

\begin{proof}
By \ref{p:relata}, the occupants of physical roles have intrinsic natures. By \ref{p:struct}, these natures are not among the structural properties that physics characterizes. By Assumption~\ref{a:phys}, every instantiated property is a fundamental property of the actual fundamental theory or is wholly grounded in such properties. The intrinsic natures of fundamental entities are not wholly grounded in structural properties alone, since by \ref{p:closure} they would then be structural, contrary to Definition~\ref{d:nature}. So either they are fundamental properties of the theory or they are grounded, at least in part, in fundamental non-structural properties; in either case there are fundamental non-structural properties, and since a fundamental property of the theory that is not structural does not itself play a dynamical role, by Assumption~\ref{a:phys} it is a categorical realizer of a property that does.
\end{proof}

Lemma~\ref{l:natures} is where the argument departs from the principal one. MND-01-a infers that physical roles are occupied by natures relevant to consciousness from the failure of the zombie hypothesis to be secondarily possible. Here the existence of non-structural fundamental natures is inferred from what it is to be a concrete relatum, and their physical status from the definition of the physical that MND-02 defends, which admits categorical realizers precisely so that physicalism would not be falsified by the structural character of physics (MND-02-a, discussion of Definition~2.2(b)). The lemma says nothing yet about consciousness.

\begin{lemma}\label{l:notfund}
No phenomenal property is fundamental. Every instantiated phenomenal property is wholly grounded in fundamental physical properties.
\end{lemma}

\begin{proof}
A phenomenal property is mental: its instantiation consists in some subject's experiencing. By Assumption~\ref{a:phys}, no fundamental property of the actual fundamental theory is mental, so no phenomenal property is fundamental, and every instantiated phenomenal property, being physical, is wholly grounded in fundamental properties of that theory.
\end{proof}

\begin{lemma}\label{l:nonstructground}
The ground of every instantiated phenomenal property includes fundamental intrinsic natures.
\end{lemma}

\begin{proof}
Let $F$ be an instantiated phenomenal property. By Lemma~\ref{l:notfund}, $F$ is wholly grounded in fundamental properties. If all of them were structural, $F$ would be structural by \ref{p:closure}, contrary to \ref{p:nonstruct}. So at least one of them is a fundamental non-structural property, which by Lemma~\ref{l:natures} is an intrinsic nature of fundamental entities.
\end{proof}

\begin{lemma}\label{l:proto}
The intrinsic natures of fundamental entities are protophenomenal.
\end{lemma}

\begin{proof}
We check the three clauses of the definition of a protophenomenal property (MND-18-a, Definition~2.3). First, the natures are not experiences: experiences are instances of mental properties, and by Assumption~\ref{a:phys} no fundamental property is mental. Second, they do not by themselves make their bearers subjects: a subject is an entity that has experiences (MND-18-a, Definition~2.1), and a fundamental entity whose fundamental properties include no mental property, and whose experience, if it had any, would require a phenomenal property grounded in a combination of natures and structure, has none by virtue of its nature alone. Third, they can in suitable combination ground experiences: by Lemma~\ref{l:nonstructground}, the natures of some fundamental entities enter into the grounds of experience in the brains of conscious subjects; by Assumption~\ref{a:unif}, every fundamental entity of a kind found in such brains has a nature of the same kind, and so a nature that, in a suitable combination, grounds experience.
\end{proof}

The third step uses Assumption~\ref{a:unif} only for the kinds of fundamental entity found in brains. If there are fundamental kinds that never occur in the brains of conscious subjects, the lemma, as proved, does not cover them; we return to this in Section~\ref{s:not}.

\begin{theorem}\label{t:main}
Given \ref{p:struct}--\ref{p:closure} and Assumptions~\ref{a:phys} and~\ref{a:unif}, protophenomenal Russellian monism is true.
\end{theorem}

\begin{proof}
Clause (i) of Definition~\ref{d:prm} is \ref{p:struct} together with Lemma~\ref{l:natures}. Clause (ii) is Lemma~\ref{l:proto}. For clause (iii): by Lemma~\ref{l:notfund} every experience is wholly grounded in fundamental properties, and by \ref{p:struct} and Lemma~\ref{l:natures} the fundamental properties are structural properties and the intrinsic natures that realize them; by Lemma~\ref{l:nonstructground} the natures are among the grounds. So experiences are wholly grounded in intrinsic natures together with the structures in which their bearers stand.
\end{proof}

Theorem~\ref{t:main} is stronger than Corollary~4.3 of MND-01-a in one respect and weaker in another. It is stronger because it yields the protophenomenal version directly, without the help of MND-18: the protophenomenal status of fundamental natures follows from physicalism, as MND-01-d notes in the first of its two routes to the same conclusion (MND-01-d, Proposition~1.1). It is weaker because it rests on \ref{p:relata}, \ref{p:nonstruct}, \ref{p:closure} and Assumption~\ref{a:unif}, which the principal argument does not need. Section~\ref{s:compare} compares the two sets of premises.

\begin{corollary}\label{c:intrinsic}
Consciousness is the intrinsic nature of physical structure: the phenomenal character of a subject's experience is a non-structural property of the physical arrangement that grounds it, wholly grounded in the intrinsic natures of the arrangement's constituents and their structure.
\end{corollary}

\begin{proof}
By Theorem~\ref{t:main}, clause (iii), the phenomenal character of an experience is wholly grounded in the intrinsic natures and structure of some arrangement of fundamental entities. By \ref{p:nonstruct}, it is not structural; by MND-10-a, Theorem~5.4, it is intrinsic to the experience. It is therefore a non-structural property of the arrangement, grounded in what the arrangement is in itself and not in its relations to other things: an intrinsic nature of that physical structure.
\end{proof}

Corollary~\ref{c:intrinsic} is the first half of the family thesis. The phrase \enquote{intrinsic nature of physical structure} is to be read with Theorem~\ref{t:main} in view: the intrinsic natures of the fundamental constituents are protophenomenal, and consciousness is the intrinsic character of certain arrangements of them, namely those that ground subjects. Which arrangements these are is the boundary problem, and which qualities they ground is the palette problem, both of which Volume~IV of MND-01 treats and MND-01-d treats in a preliminary way.

We now derive the second half, that the explanatory gap arises from two modes of presentation. For this we need one epistemic premise, which MND-01-b states and defends: whatever is a priori derivable from truths stated in the structural mode is itself statable in the structural mode (MND-01-b, P2).

\begin{corollary}\label{c:gap}
Let $P$ be the conjunction of microphysical truths and $Q$ a phenomenal truth. Given Theorem~\ref{t:main} and the epistemic principle of structural closure, the conditional $P \to Q$ is necessary and not a priori, and the explanatory gap arises from the difference between the structural mode in which physics presents the arrangements that ground experience and the intrinsic mode in which experience presents them.
\end{corollary}

\begin{proof}
Necessity: by Theorem~\ref{t:main}, the phenomenal property that $Q$ attributes is wholly grounded in the intrinsic natures and structure of an arrangement of the entities that $P$ describes; read as concerning the actual occupants of the roles it describes, $P$ concerns that arrangement, and full grounding entails necessitation. Failure of apriority: by \ref{p:struct}, $P$ is stated in the structural mode; by \ref{p:nonstruct}, $Q$ is not statable in it; so by the epistemic principle of structural closure, $Q$ is not a priori derivable from $P$. The rest is Theorem~4.3 of MND-01-b, whose premises are now all available.
\end{proof}

Corollary~\ref{c:gap} reproduces Corollary~4.4 of MND-01-a, that given materialism the zombie hypothesis is not secondarily possible and is not ruled out a priori, without any premise about conceivability. In particular, it shows that zombies are conceivable in the only sense the gap requires, that their absence is not a priori, and impossible, as MND-05 concludes (MND-05-a, Theorem~3.3). The conceivability of zombies is here a consequence of the view, not a premise of the argument for it.

\begin{proposition}[Completeness of the completed picture]\label{p:complete}
Given the premises of Theorem~\ref{t:main} and the completeness of grounding (MND-02-a, Assumption~3.4), every instantiated property is of exactly one of three kinds: a structural property, an intrinsic nature of fundamental entities, or a property wholly grounded in properties of the first two kinds and not itself of either. Phenomenal properties are of the third kind, and are the only known properties of that kind whose instances we can identify a priori as not structural.
\end{proposition}

\begin{proof}
By Assumption~\ref{a:phys} and MND-02-a, Assumption~3.4, every instantiated property is fundamental or wholly grounded in fundamental properties; by \ref{p:struct} and Lemma~\ref{l:natures}, fundamental properties are structural or intrinsic natures. Phenomenal properties are not fundamental (Lemma~\ref{l:notfund}) and not structural (\ref{p:nonstruct}), so they are of the third kind. That their non-structural status is knowable by the subject is \ref{p:nonstruct} as supported by MND-10; for the other non-fundamental properties we know of, we have either a structural characterization or no characterization at all.
\end{proof}

Proposition~\ref{p:complete} is the sense in which the view completes physics. Physics as we have it describes the first kind of property. The second kind is required by \ref{p:relata} and admitted by the definition of the physical of MND-02. The third kind is everything else, and among it are some properties, the phenomenal ones, that are not structural and that we know in the intrinsic mode. There is no further kind: in particular, there is no fundamental mental property and no property that floats free of the grounding structure. The completed picture has no place for the hard problem as an ontological problem; it retains it as an epistemic one, located by Corollary~\ref{c:gap}.

\section{The causal role of experience}\label{s:causal}

A further advantage of the completed picture is that it makes room for the causal efficacy of experience. We do not use this in the argument for Theorem~\ref{t:main}; it is a consequence, and a test the view must pass.

The exclusion argument concludes that if mental properties are not identical to physical ones, they are causally idle, since every physical effect has a sufficient physical cause \citep{Kim2005}. On a dualism of fundamental phenomenal properties, the argument has force: such properties would either overdetermine physical effects or be epiphenomenal. MND-11 argues that mental properties can be causes because they are proportional to their effects, whereas their realizers are overly specific (MND-11-a, Theorem~3.4). The completed picture is consistent with that result and adds something to it.

\begin{proposition}\label{p:causal}
On protophenomenal Russellian monism, the phenomenal character of an experience is wholly grounded in properties that are causally efficacious in the physical world, and is not an additional cause of any physical effect.
\end{proposition}

\begin{proof}
By Theorem~\ref{t:main}, the phenomenal character is wholly grounded in intrinsic natures and structure. By Lemma~\ref{l:natures}, the natures are the categorical realizers of the properties that play roles in the dynamical laws, and the structure is the pattern of those roles. Both are therefore involved in physical causation: the natures are what does the work that the structural properties describe. Since the phenomenal character is wholly grounded in them, it adds nothing to what there is, and so is not an additional, overdetermining cause.
\end{proof}

Proposition~\ref{p:causal} shows that experience, on the completed picture, is neither epiphenomenal nor an overdeterminer. Whether it is a cause in its own right, at its own level, is the question MND-11 answers by proportionality: where a phenomenal property is the proportional cause of a behavioral effect, as when the painfulness of a stimulus, and not the exact neural configuration that realizes it, makes the difference to withdrawal, it is the cause at the grain of the effect. The completed picture adds that the realizers which proportionality treats as overly specific are not themselves purely structural: their own causal work is done by intrinsic natures, and experience is grounded in those natures. The physical causation that the exclusion argument sets against mental causation is therefore not a purely structural causation from which experience is absent. This is the reason for which \citet{Chalmers1996} counted Russellian monism as the view that best reconciles the anti-materialist arguments with causal closure, and it is available on the present route without those arguments.

\section{The two routes compared}\label{s:compare}

The two arguments for the family's conclusion can be compared by their premises.

\begin{center}\small
\begin{tabular}{@{}p{0.42\textwidth}cc@{}}
\toprule
Premise & MND-01-a & MND-01-c \\
\midrule
Structuralism about physics (MND-01-a, P2) & used & used \\
Materialism (MND-02-a, Theorem~3.6) & used & used \\
Zombies ideally conceivable (4.1(a)) & used & not used \\
Ideal conceivability to primary possibility (4.1(b)) & used & not used \\
Primary possibility to secondary or Russellian (4.1(c)) & used & not used \\
Secondary possibility refutes materialism (4.1(d)) & used & not used \\
Relata have intrinsic natures (\ref{p:relata}) & not used & used \\
Phenomenal properties not structural (\ref{p:nonstruct}) & implicit & used \\
Structural grounding closure (\ref{p:closure}) & not used & used \\
Uniformity of natures (Assumption~\ref{a:unif}) & not used & used \\
\bottomrule
\end{tabular}
\end{center}

Three observations follow. First, the routes share their non-modal core, structuralism and materialism, and differ in what they add. The principal route adds modal premises, connecting what can be conceived with what is possible in two senses. The present route adds metaphysical premises about relations, grounding and the uniformity of nature. A critic who rejects the conceivability--possibility link, as type-B physicalists do, must accept the second route if she accepts its metaphysical premises; a critic who rejects the metaphysical premises, as a pure-powers ontologist might, must accept the first route if she accepts its modal premises. The family conclusion is therefore supported by two arguments with disjoint contested premises.

Second, the non-structurality of phenomenal properties is implicit in the principal argument: clause (c) of Assumption~4.1 is plausible only because phenomenal terms have the same primary and secondary intension, which is to say that they pick out phenomenal properties by what they are like rather than by a role (MND-01-a, discussion of Assumption~4.1). The present route makes the premise explicit and rests it on MND-10's argument from Newman's point rather than on a claim about intensions.

Third, the present route has a dialectical advantage against the type-B physicalist. She holds that phenomenal properties are identical with physical properties and that the gap is due to the isolation of phenomenal concepts. Theorem~\ref{t:main} shows that if she also accepts \ref{p:struct}--\ref{p:closure}, the physical properties with which phenomenal properties are identical cannot be purely structural, and so must involve intrinsic natures. Type-B physicalism, so constrained, becomes a form of Russellian monism. The constraint is not imposed from outside: \ref{p:struct} is a claim about physics that most type-B physicalists accept, and \ref{p:nonstruct} is the claim, which they also accept, that phenomenal concepts do not present their referents as role-occupants \citep{Loar1990}, together with the claim that what they present is how the referents are.

\section{Objections and replies}\label{s:obj}

\begin{objection}[Pure powers]\label{o:powers}
\ref{p:relata} is false if fundamental properties are pure powers: properties whose whole nature is dispositional, and which therefore need no categorical nature behind them \citep{Bird2007}. MET-11 itself holds that dispositions are irreducible to categorical bases.
\end{objection}

\begin{reply}
MET-11 argues that dispositions are real and irreducible properties of their bearers, and in particular that they are not reducible to categorical bases together with laws (MET-11-a, Theorem~3.4). \ref{p:relata} is consistent with this. It does not say that dispositions reduce to or are grounded in intrinsic natures; it says that the bearers of dispositions have intrinsic natures in addition. A view on which each fundamental property is both irreducibly dispositional and qualitative, the identity theory of powerful qualities \citep{Heil2003}, satisfies \ref{p:relata} and is consistent with MET-11. What \ref{p:relata} excludes is the stronger view on which the bearers of dispositions have \emph{no} properties except dispositional ones, so that the world is relations all the way down. Against that view the regress argument applies: each power is individuated by its manifestations, which are the instantiation of further powers, individuated by theirs, and nothing at any point is specified except a place in a pattern. The pure-powers theorist has answers to the charge that this is a vicious regress \citep{Bird2007}, but none that shows how such a world would differ from a mathematical structure of the same cardinality. The present argument does not need to refute the pure-powers view to be a proof of the family's conclusion from its premises; it needs only that \ref{p:relata} is true, and the argument from concreteness supports it.
\end{reply}

\begin{objection}[Humility]\label{o:humility}
If physics tells us only about roles, we are irremediably ignorant of the intrinsic natures that occupy them \citep{Langton1998,Lewis2009}. The completed picture claims to know that those natures are protophenomenal, and so claims a knowledge that structuralism itself rules out.
\end{objection}

\begin{reply}
Humility, as Lewis defends it, is the thesis that we cannot know which intrinsic properties occupy which roles, since any permutation of the occupants that preserves the roles is consistent with all our evidence. Theorem~\ref{t:main} is consistent with this. It does not say which nature occupies which role, or what any individual nature is like. It says that the natures, whatever they are, are not mental and can ground experience in combination, and it reaches this by a general argument, not by identifying any of them. The argument uses knowledge of a different kind: knowledge that experience exists and is not structural, which comes from introspection (MND-10-a, P1), and knowledge that everything is physical, which comes from physics and physiology (MND-02-a). Humility forbids knowing the occupants one by one, not knowing general truths about them from premises of this kind. MND-01-b argues in addition that we have partial access to the intrinsic character of the arrangements that ground our own experiences, by acquaintance; that too is consistent with humility about the fundamental natures taken singly.
\end{reply}

\begin{objection}[Emergence]\label{o:emergence}
\ref{p:closure} begs the question against the emergentist, who holds that experience emerges from purely structural organization, as a novel property that is not itself structural.
\end{objection}

\begin{reply}
Two kinds of emergence must be distinguished. On the first, the emergent property is wholly grounded in its base. Then it is wholly grounded in structure alone, and the question is whether a property so grounded can be non-structural. \ref{p:closure} says it cannot, because a property whose instantiation is wholly grounded in the instantiation of relations has instantiation conditions consisting in the instantiation of relations. The emergentist must deny this and owes an account of how a ground that consists in nothing but relations can make it the case that something is non-relationally some way; the hard problem is, in one of its forms, the demand for such an account, and its persistence is evidence that none is available. On the second kind of emergence, the emergent property is not wholly grounded in its base but is connected to it by a fundamental law. Then it is an ungrounded mental property, and by Proposition~3.5 of MND-02-a it is not physical. Emergence of the second kind is therefore inconsistent with Assumption~\ref{a:phys}, which the emergentist of this kind rejects. The argument is not addressed to her; she is addressed by the case for physicalism in MND-02.
\end{reply}

\begin{objection}[Illusionism]\label{o:illusion}
\ref{p:nonstruct} is denied by illusionists, who hold that there are no phenomenal properties as \ref{p:nonstruct} conceives them, and by strong representationalists, who hold that phenomenal character is exhausted by representational content, which is relational.
\end{objection}

\begin{reply}
Both positions are addressed in MND-10, whose arguments we take over: illusionism is false because the seeming it posits is either itself phenomenal or fails to explain the datum (MND-10-a, Lemma~4.2), and qualia are not constituted by the relations of experiences to what they represent (MND-10-a, Lemma~5.3). The principal route also faces the illusionist, who rejects Assumption~4.1(a). The two routes thus share one opponent here; neither is weaker than the other against her.
\end{reply}

\begin{objection}[Idle natures]\label{o:idle}
Protophenomenal natures are posited to ground experience, but they are said to do nothing that the structural properties do not already describe. A posit that adds nothing to the explanation of any physical event is idle, and should be eliminated.
\end{objection}

\begin{reply}
The natures are not idle in physics: by Lemma~\ref{l:natures} they are the categorical realizers of the properties that play roles in the laws, and so they are what the laws describe the behavior of. They add nothing to the \emph{description} that physics gives, since that description is structural, but they are required by the existence of what it describes, by \ref{p:relata}. Nor are they idle in the theory of consciousness: by Lemma~\ref{l:nonstructground} they are indispensable to the ground of experience. The demand that every posit add to the structural description of the world is the demand that the world be nothing but structure, which \ref{p:relata} denies.
\end{reply}

\section{What the argument does not establish}\label{s:not}

The argument establishes protophenomenal Russellian monism for the kinds of fundamental entity that compose the brains of conscious subjects. It does not establish anything about the intrinsic natures of kinds, if there are any, that never occur in such brains, such as particles that interact only gravitationally. \ref{p:relata} gives them intrinsic natures, and Lemma~\ref{l:natures} makes those natures physical, but Lemma~\ref{l:proto} does not show that they are protophenomenal. Nothing in the family's conclusion requires that it should.

The argument does not solve the combination problem in its residual forms. It shows that experience is grounded in protophenomenal natures and structure, and therefore that the grounding of experience in micro-experiences, which MND-18 shows to be impossible, is not required (MND-18-a, Proposition~4.1). It does not say which arrangements ground subjects or which qualities they ground. These are the boundary and palette problems, which MND-01-d treats as problems of articulating the grounding principles and which Volume~IV of MND-01 develops; on the treatment proposed there, the boundaries of a subject are the boundaries of a maximal integrated arrangement.

The argument does not establish physicalism. It takes the physicalism of MND-02 as a premise, as does the principal argument. Without it, Lemma~\ref{l:notfund} fails, and the remaining premises are consistent with a dualism of fundamental phenomenal properties and with panpsychism.

Finally, the argument does not show that the principal argument is unsound. Its purpose is to show that the family's conclusion does not depend on Assumption~4.1 of MND-01-a, and so survives the rejection of any of its clauses.

\section{Conclusion}\label{s:concl}

The conclusion of MND-01 can be reached without the two-dimensional argument. Physics describes the world structurally; relations require relata with intrinsic natures; phenomenal properties are not structural; everything is physical in the sense of MND-02; and structure alone grounds only structure. It follows that the fundamental entities of physics have intrinsic natures that are physical and not mental, that experience is wholly grounded in those natures together with structure, and that the natures are protophenomenal. Consciousness is then the intrinsic nature of certain physical structures, and the explanatory gap is the gap between the structural mode in which physics presents them and the intrinsic mode in which experience does. The argument avoids every clause of Assumption~4.1 of MND-01-a, and in particular the inference from ideal conceivability to primary possibility.

\begin{thebibliography}{99}\small
\bibitem[Alter and Nagasawa(2012)]{AlterNagasawa2012} Alter, T. and Nagasawa, Y. (2012). What is Russellian monism? \emph{Journal of Consciousness Studies}, 19(9--10), 67--95.
\bibitem[Alter and Nagasawa(2015)]{AlterNagasawa2015} Alter, T. and Nagasawa, Y. (Eds.) (2015). \emph{Consciousness in the Physical World: Perspectives on Russellian Monism}. Oxford University Press.
\bibitem[Bird(2007)]{Bird2007} Bird, A. (2007). \emph{Nature's Metaphysics: Laws and Properties}. Oxford University Press.
\bibitem[Blackburn(1990)]{Blackburn1990} Blackburn, S. (1990). Filling in space. \emph{Analysis}, 50(2), 62--65.
\bibitem[Block and Stalnaker(1999)]{BlockStalnaker1999} Block, N. and Stalnaker, R. (1999). Conceptual analysis, dualism, and the explanatory gap. \emph{The Philosophical Review}, 108(1), 1--46.
\bibitem[Chalmers(1996)]{Chalmers1996} Chalmers, D.~J. (1996). \emph{The Conscious Mind: In Search of a Fundamental Theory}. Oxford University Press.
\bibitem[Chalmers(2009)]{Chalmers2009} Chalmers, D.~J. (2009). The two-dimensional argument against materialism. In B.~McLaughlin, A.~Beckermann and S.~Walter (Eds.), \emph{The Oxford Handbook of Philosophy of Mind} (pp.~313--335). Oxford University Press.
\bibitem[Chalmers(2017)]{Chalmers2017} Chalmers, D.~J. (2017). The combination problem for panpsychism. In G.~Br\"untrup and L.~Jaskolla (Eds.), \emph{Panpsychism: Contemporary Perspectives}. Oxford University Press.
\bibitem[Coleman(2014)]{Coleman2014} Coleman, S. (2014). The real combination problem: panpsychism, micro-subjects, and emergence. \emph{Erkenntnis}, 79(1), 19--44.
\bibitem[Demopoulos and Friedman(1985)]{DemopoulosFriedman1985} Demopoulos, W. and Friedman, M. (1985). Bertrand Russell's \emph{The Analysis of Matter}: its historical context and contemporary interest. \emph{Philosophy of Science}, 52(4), 621--639.
\bibitem[Eddington(1928)]{Eddington1928} Eddington, A.~S. (1928). \emph{The Nature of the Physical World}. Cambridge University Press.
\bibitem[Feigl(1958)]{Feigl1958} Feigl, H. (1958). The \enquote{mental} and the \enquote{physical}. In H.~Feigl, M.~Scriven and G.~Maxwell (Eds.), \emph{Minnesota Studies in the Philosophy of Science}, Vol.~2 (pp.~370--497). University of Minnesota Press.
\bibitem[Goff(2017)]{Goff2017} Goff, P. (2017). \emph{Consciousness and Fundamental Reality}. Oxford University Press.
\bibitem[Heil(2003)]{Heil2003} Heil, J. (2003). \emph{From an Ontological Point of View}. Oxford University Press.
\bibitem[Kim(2005)]{Kim2005} Kim, J. (2005). \emph{Physicalism, or Something Near Enough}. Princeton University Press.
\bibitem[Langton(1998)]{Langton1998} Langton, R. (1998). \emph{Kantian Humility: Our Ignorance of Things in Themselves}. Oxford University Press.
\bibitem[Lewis(2009)]{Lewis2009} Lewis, D. (2009). Ramseyan humility. In D.~Braddon-Mitchell and R.~Nola (Eds.), \emph{Conceptual Analysis and Philosophical Naturalism}. MIT Press.
\bibitem[Loar(1990)]{Loar1990} Loar, B. (1990). Phenomenal states. \emph{Philosophical Perspectives}, 4, 81--108.
\bibitem[Lockwood(1989)]{Lockwood1989} Lockwood, M. (1989). \emph{Mind, Brain and the Quantum: The Compound \enquote{I}}. Blackwell.
\bibitem[Maxwell(1978)]{Maxwell1978} Maxwell, G. (1978). Rigid designators and mind--brain identity. In C.~W. Savage (Ed.), \emph{Perception and Cognition: Issues in the Foundations of Psychology}, Minnesota Studies in the Philosophy of Science, Vol.~9. University of Minnesota Press.
\bibitem[Newman(1928)]{Newman1928} Newman, M.~H.~A. (1928). Mr.\ Russell's \enquote{causal theory of perception}. \emph{Mind}, 37(146), 137--148.
\bibitem[Papineau(2001)]{Papineau2001} Papineau, D. (2001). The rise of physicalism. In C.~Gillett and B.~Loewer (Eds.), \emph{Physicalism and Its Discontents}. Cambridge University Press.
\bibitem[Papineau(2002)]{Papineau2002} Papineau, D. (2002). \emph{Thinking about Consciousness}. Oxford University Press.
\bibitem[Pereboom(2011)]{Pereboom2011} Pereboom, D. (2011). \emph{Consciousness and the Prospects of Physicalism}. Oxford University Press.
\bibitem[Russell(1927)]{Russell1927} Russell, B. (1927). \emph{The Analysis of Matter}. Kegan Paul.
\bibitem[Stoljar(2001)]{Stoljar2001} Stoljar, D. (2001). Two conceptions of the physical. \emph{Philosophy and Phenomenological Research}, 62(2), 253--281.
\bibitem[Strawson(2006)]{Strawson2006} Strawson, G. (2006). Realistic monism: why physicalism entails panpsychism. \emph{Journal of Consciousness Studies}, 13(10--11), 3--31.
\end{thebibliography}

\end{document}
